\clearpage
\section{Einleitung}

Das Projekt im Vertiefungsprofil GeoSensorik \& Monitoring II ist in die Teilprojekte A (Literaturrecherche), 
B (Auswertung Schwanderbärgli) und C (Geomonitoring-Projekt am Campus) aufgeteilt.
Die Teilprojekte werden unabhängig voneinander bearbeitet und im vorliegenden Bericht dokumentiert. 

Teilprojekt A enthält eine umfassende Literaturrecherche zum Thema 'Zenith Total Delay' im Bezug auf GNSS-Messungen.
Dabei wird weiter auf die Anwendung in der Meteorologie und den Nutzen von Mapping Functions eingegangen. 

Teilprojekt B umfasst die Auswertung und Dokumentation der im Feldkurs 2026 erfolgten GNSS-Messungen,
sowie den Vergleich mit der Trigonometrischen Höhenbestimmung und die Auswertung der Scanning-Daten.
Die übrigen Auswertungen des Schwanderbärgli sind in den anderen Berichten (der Teilgruppe 2 bis 4) dokumentiert.

Teilprojekt C beinhaltet das Geomonitoring-Projekt am Campus Muttenz in welchem der Aufbau eines
einfachen Monitoring-Systems mittels Raspberry-Pi und Tachymeter untersucht wurde.
Im Weiteren werden die erfassten Daten in einer Zeitreihenanalyse ausgewertet und interpretiert.

Das abschliessende Fazit ist ebenfalls unterteilt in die behandelten Teilprojekte.